\chapter{Conclusion generale}

\section{Introduction du chapitre}
Ce dernier chapitre clot le rapport en synthetisant les apports principaux du projet et leur impact sur la securisation des acces privilegies dans une logique Zero Trust.

Ce projet a permis de concevoir et de valider une architecture Zero Trust adaptee a un besoin concret de maitrise des acces privilegies.

La contribution principale est la mise en place d'un parcours complet, depuis l'authentification forte jusqu'a la revocation automatique des droits, en passant par l'approbation humaine et la tracabilite.

Les objectifs initiaux ont ete couverts:
\begin{itemize}
	\item centraliser la verification d'identite via IAM;
	\item controler l'acces a la ressource via ACL et workflow d'approbation;
	\item rendre l'elevation de privilege temporaire, justifiee et auditable.
\end{itemize}

Au-dela de la realisation technique, le projet montre qu'une demarche Zero Trust est d'abord une logique d'ingenierie et de gouvernance: chaque acces doit etre explicite, limite dans le temps, et verifiable a posteriori.

\section{Retour sur la problematique}
La problematique initiale consistait a supprimer les acces permanents tout en conservant l'operationnel. Le prototype apporte une reponse pragmatique grace au JIT, a la journalisation et a l'approbation explicite.

\section{Contributions}
Les contributions principales sont la reduction de la surface d'attaque, l'amelioration de la tracabilite et la structuration d'un workflow exploitable pour la conformite.

\section{Ouverture}
Les evolutions futures peuvent aller vers un PAM plus complet, davantage d'automatisation et une meilleure integration avec les outils de ticketing et de politique centralisee.

\section{Conclusion du chapitre}
Le projet atteint son objectif central: transformer l'acces privilegie permanent en acces conditionnel, temporaire et auditable. Cette base constitue un socle solide pour une generalisation progressive a d'autres environnements et usages.
